\cotefile{156-4}
\scanpage{66}\mspage{65}
Ceci résulte du \deleted{\uline{C}4 et} cor. 1 et du

\uline{Lemme} Pour que $F$ soit de t. f.\inserted{(resp. \unsure{fini|finie})} Il f et s que a) les $F_i$ le soient et b) Tous les $F_i$ sauf un \unsure{nb} fini sont $\emptyset$.

Notations $\pi_0(F)$ (NB Si $F\leq G$ on a $\pi_0(F)\longrightarrow\pi_0(G)$\deleted{\unsure{)|?}}, et transitivité) $\pi_0(M)$, $\pi_0^{\deleted{\unsure{\mathrm{cat}|\mathrm{act}}}!}(\mathcal F)$. (\unsure{Attention}, le $\pi_0^{\unsure{d|!}}(\mathcal F)$ n'est pas le $\pi_0(F,\leq)$, mais le $\pi_0$ ``en tenant compte \unsure{du|des} \unsure{Sup}''. Si $\unsure{\mathcal C|\mathscr C|C}=(\mathcal F,\leq,\deleted{\unsure{\leq|\prec}}\unsure{\preceq|\leqslant})$\margin{\unsure{Contrée|Contrie}}\margin{qui ne sert \deleted{\unsure{pourtant}} \unsure{à} rien ici}, on pose aussi $\pi_0(\mathcal C)$.

Axiome renforcé

\uuline{C4 bis} Pour tout $(F_i)_i\in\prod\mathcal F_i$, $\deleted{\unsure{\exists F|\mathcal F}}$ $\exists F\in\mathcal F$ telle que $F\mid M_i=\deleted{\unsure{\mathcal F}}F_i$ pour $i\unsure{\vdash|\in}I$.

\uline{\unsure{Contrée|Contrie} \unsure{résiduelle|réssiduelle}} \uline{Support} \uline{et} \uline{cosupport}

Soit $\deleted{\unsure{(F_i)\in\mathcal F|(H_i)\in\mathcal F}}$ $\Phi\subset\mathcal F$ un ens. de figures \inserted{$H_i$}. \deleted{On posera} Soit $F$ une figure. Pour \deleted{\unsure{Cosupp.} $\Phi$} que $F$ soit \unsure{disjointe|disjoint} de tous les $H_i$, il f et s. que ce soit le cas pour toutes les strates de $F$ (à cause de C3). \deleted{\unsure{Posons|Poser}} Posons
\scanpage{67}\mspage{66}
\[
\begin{aligned}
\mathcal F'&=\unsure{\operatorname{cosup}|\operatorname{cosupp}}_{\mathcal F}(\Phi)=\{\,F\in\mathcal F\mid F\ \text{disjoint \unsure{des|de} }H_i\unsure{\vdash|\in}\Phi\inserted{\subset\mathcal F}\,\}\\
M'&=\operatorname{cosup}_{M}(\Phi)=\{\,X\in M\mid X\ \text{disjoint \unsure{de|des} }H_i\in\unsure{\Phi|\varphi}\unsure{)|\}}\,\}\\
&\inserted{=M\cap\operatorname{cosup}_{\mathcal F}(\Phi)\subset M}
\end{aligned}
\]

\deleted{Donc \unsure{on aura|on a}, pour $F\in\mathcal F$ posant $\mathcal F'=\operatorname{cosup}_{\mathcal F}$,}

\[
\begin{aligned}
&\deleted{\unsure{\phi|\Phi}}\qquad
\deleted{\mathcal F'=\operatorname{cosupp}_{\mathcal F}(\Phi),\quad M'=\mathcal F'\cap M,}\\[-2pt]
&\qquad\deleted{F\in\operatorname{cosup}_{\unsure{,|\mathcal F}}}
\qquad\qquad\deleted{=\operatorname{cosup}_{\deleted{\unsure{\mathcal F}}M}(\mathcal F)}
\end{aligned}
\]

Donc pour $F\in\mathcal F$
% slanted
\margin{Ainsi, $M'$ et $\mathcal F'$ se déterminent \unsure{mutuellement}}

\[
\left\{\begin{aligned}
F\in\mathcal F'&\Longleftrightarrow\overline F\subset M'\\
M'&=\unsure{.|\cdot}\mathcal F'\cap M
\end{aligned}\right.
\]
\margin{NB $\mathcal F',M'$ sont des parties fermées de $\mathcal F,M$. $\mathcal F'$ \uline{\unsure{Stable|Stables} par Sup.}}

Soit d'autre part
\[
\begin{aligned}
\mathcal F''&=\operatorname{cosupp}_{\mathcal F}(\mathcal F')\overset{\text{def}}{=}\operatorname{supp}_{\mathcal F}(\Phi)\\
M''&=\operatorname{cosup}_{M}(\mathcal F')\overset{\text{def}}{=}\operatorname{supp}_{M}(\Phi)
\end{aligned}
\]

On a \unsure{bien|vu|la|là} $\mathcal F''\supset\Phi$, $\mathcal F''$ stable par \unsure{Sup|Supp}, \deleted{contient $\Phi$} et par passage \unsure{à|:} une sous-figure. Soit $\deleted{\unsure{M}} \overline\Phi=\bigcup_{H_i\in\Phi}\overline H_i$, alors
\[
\begin{aligned}
\operatorname{cosup}_{\mathcal F}(\Phi)&=\operatorname{cosup}_{\mathcal F}(M_{\overline\Phi})\\
\operatorname{cosup}_{\deleted{\unsure{\mathcal F}}M}(\Phi)&=\{\,\deleted{\unsure{\varnothing|x}}X\in M\mid X\ \text{disjoint de }\unsure{\mathrm{tt}|\forall}\ Y\in M_{\overline\Phi}\,\}
\end{aligned}
\]
donc on est \unsure{ramené|revenu} à des constructions dans $M$. \unsure{\S|?}

\uline{Partie saturée \inserted{$\mathcal F_0$} de} $\deleted{\unsure{M}}\mathcal F$ : égale au cosupport de son cosupport ($\text{p. r. à }\mathcal F$). Partie \unsure{saturée} \inserted{$M_0$} de $M$ : égale au cosupport de son support (dans $M$). Elles se correspondent \unsure{mutuellement} \unsure{par}
\[
\begin{aligned}
M_0&=\mathcal F_0\cap M\\
\mathcal F_0&=\{\,F\in\mathcal F\mid\overline F\subset M_0\,\}
\end{aligned}
\]
\scanpage{68}\mspage{67}
\deleted{On a une notion de \unsure{``c|``C|``6}}

Toute partie de $\mathcal F$ (resp $M$) est contenue dans une partie saturée \unsure{minimale} (parmi \unsure{celles|telles} la contenant), qui \deleted{s'appelle \unsure{du|de}} n'est autre que le cosupport du cosupport, et s'appelle aussi son support. Les parties saturées sont celles qui sont égales au support, ou au cosupport, d'une partie. L'application \unsure{.} $\unsure{\operatorname{cosup}|\operatorname{cosupp}}$ définit une \unsure{bij\inserted{e}ction} de l'ens des parties \unsure{saturées|saturée} \unsure{sur} elles \unsure{=|mêmes}, renversant la relation d'inclusion.

On désigne par $\Sigma$ l'ens des \uline{supports} \uline{virtuels}, ou parties saturées \inserted{$M_0$} de $M$, \unsure{où|ou} \deleted{\unsure{de}} $\mathcal F_0$ de $M$. \deleted{On \unsure{définit}} \deleted{L'application du cosupport de} On \unsure{désidentifie|délibérément identifie|identifie} $\Sigma$ de \deleted{$\unsure{\mathcal F}$} ses deux \unsure{réalisations} (comme \unsure{partie|parties} de $\unsure{\mathscr P|\mathcal P|\wp}(M)$, ou de $\mathscr P(\mathcal F)$) et on écrit, pour $\sigma\in\Sigma$
\[
M_\sigma\subset M,\qquad\mathcal F_\sigma\subset\mathcal F
\]

Dans $\Sigma$, \deleted{pour} la notion de cosupport est une anti\deleted{\unsure{isomorph|homéo}}\unsure{involution|bijection|dualité} d'ens. ordonnés. Dans $\Sigma$, on a des $\mathrm{Inf}$ quelconques (c'est l'intersection des parties de $M$, ou de $\mathcal F$), et des $\mathrm{Sup}$ quelconques (c'est la partie saturée
\scanpage{69}\mspage{68}
engendrée par la réunion) Je \deleted{dis} \inserted{présume} qu'on a en fait une algèbre de Boole (\inserted{``ultra}stonienne'', à cause des $\mathrm{Sup}$ \unsure{quelc.})
% slanted
\margin{Cela sera valable dans ``le cas ensembliste'', et en tous cas dans les exemples types, donnés \deleted{\unsure{tout}} au début, en termes d'un $E$ (top. ou mesurable), alors $\Sigma$ s'identifie à l'ens des parties de $E$, et \unsure{on trouve|on trouvera|on retrouve} le compactifié de Stone de $E_{\mathrm{disc.}}$ Dans le cas qui vient, \unsure{je} ne vois comment \unsure{prouver}, pour trois supports \[P,Q,R\quad P\wedge(Q\vee R)=\unsure{(|\{}(P\wedge Q)\vee(P\wedge R)\unsure{.|,|)}\]}

\deleted{Il faudra y revenir}

\unsure{Donc} les supports \unsure{s'identifient|s'identifieraient} aux \unsure{Parties|Partie} à la fois ouvertes et fermées du spectre, qui serait un espace \deleted{\unsure{to|tot}}\unsure{compact|cpct} \deleted{\unsure{tot}} \unsure{disconnexe} ultra\unsure{.|-}stonien.

L'involution correspond au passage au complémentaire.

On fera attention que la construction $\Sigma$ ne dépend que de l'ens ordonné $M$ \inserted{muni de la relation $\unsure{\mathcal R|R}$ de compatibilité} (sans connaître \inserted{\unsure{toute|toutes}} la famille des parties $\unsure{\overline\Phi|\Phi'|\Phi}$ de $M$, correspondant aux figures). En effet, en termes de $(M,\leq,R)$, le cosupport d'une partie $M'$ de $M$ est défini par
\[
\unsure{\operatorname{cosup}|\operatorname{cosupp}}_M(M')=\{\,x\in M\mid(x,x')\in\mathcal R\quad\forall x'\in M'\,\}
\]
% slanted
\margin{\uline{Ex} de ``supports'', ou parties saturées de $M$ : les comp. connexes de $M$ (utiliser C4).}
\scanpage{70}\mspage{69}
J'en viens à la rel d'ordre $\unsure{\preceq|\leqslant}$ (subdivision), de nature \unsure{très|trs} différente. Il n'y \unsure{aura} pas de plus petit (ni de plus grand) élément, pas de $\mathrm{Sup}$ de deux \deleted{\unsure{éléments|élémen}} figures \inserted{$F',F''\preceq-$} \unsure{en général|en position} \unsure{(i.e.|liée.|lire.} Subdivisions d'une \unsure{=|même} figure $F$, mais parfois un $\inf$.

On dit que deux subdivision $F',\unsure{F'''|F''}$ de $F$ sont \uline{comp\deleted{\unsure{ati|atib}}\unsure{ara|ati|ata}bles} \deleted{\unsure{on dit|c.à.d.} que \unsure{sub}} (je ne dirai pas ``compatibles'', vu que le terme est déjà pris) si elles admettent une subdivision commune.

\diagram{\begin{tikzcd}[ampersand replacement=\&,row sep=0.1em,column sep=0.25em]
\& F' \arrow[dr,phantom,"{\scriptstyle\rlap{\ensuremath{\deleted{\unsure{\leq|\prec}}}}\preceq}" description,sloped] \& \\
F''' \arrow[ur,phantom,"\scriptstyle\preceq" description,sloped] \arrow[dr,phantom,"\scriptstyle\preceq" description,sloped] \& \& F \\
\& F'' \arrow[ur,phantom,"\scriptstyle\preceq" description,sloped] \&
\end{tikzcd}}
% slanted
\margin{On sera amené à renforcer encore cette condition sur les contrées modérées !}

On dit que la \uuline{\unsure{contrée|contrie}} est \uuline{modérée} si deux subdivisions \deleted{\unsure{d'un}} \unsure{d'une =|d'un =} $F$ sont toujours comparables. \deleted{Dans le cas qui vient,} \deleted{on a}.

\deleted{\uline{C5}} Je ne saurais donner aucune propriété particulière de $\preceq$ \deleted{tout} seul. Les propriétés que j'ai en vue concernent ses relations avec $\leq$.

% slanted
\margin{\unsure{Reporter|Rejeter} après C6, $C7_0$ \unsure{ces|mes} axiomes de \uline{disjonction}}
\uline{C5} Soit $F'\preceq F$. Alors pour \unsure{tte} figure $G$, $G$ est disjointe de $F$ ssi elle est disjointe de $F'$. [\uline{NB} Il suffit de le poser pour $\deleted{\unsure{F}}\unsure{G|F}=\unsure{X|x}\in M$, à cause de C3]
\scanpage{71}
\mspage{\deleted{\unsure{69|79}}\unsure{70|79}}

Il revient au $=$ de dire que $F$ et $F'$ ont $=$ cosupport, ou encore, même support.

\unsure{Donnons nous} un support.

% Two short columns of marks in the left margin opposite this line.
\margin{$\unsure{!|\vdots}\qquad\unsure{!|\vdots}$}
Une autre façon de dire que $F,G\ \text{disj.}\Longrightarrow F',G\ \text{disj}$, est ainsi :

\uline{Corollaire} Le cosupport \inserted{dans $\mathcal F$} de tout objet $G$ de $\mathcal F$ est fermé, non \unsure{seult|seulemt} pour $\leq$, mais aussi pour $\unsure{\preceq|\leqslant}$ -- c'est donc ce qu'on appelle une \uline{sous}-\unsure{contrée|contrie}.

Donc \unsure{itou} pour toute partie $\Phi\subset\mathcal F$.

En d'autres termes, \deleted{toute} partie de $\mathcal F$ qui est un $\mathcal F$-support (ou ``$\mathcal F$-saturée, sous-entendu : en termes de $\leq$\unsure{,)|,) ;} cf page $67$) est \unsure{automatiquemt|automatiquement} \deleted{\unsure{f}}une sous-contrée.

% Cancelled draft, written upside down. The diagonal strokes run across the whole block.
\mspage{\deleted{\unsure{5\deleted{\unsure{4|3}}5}}}

\deleted{un ensemble, et $\deleted{\unsure{\mathcal F|F}}\Phi$ \unsure{un} ens \unsure{de|des} parties de $M$, satisfaisant la condition précédente \unsure{(au --)|(ou --)|(en --)}.}

\deleted{Alors $\Phi$, muni de la relation d'inclusion, satisfait $C\,1$, et ses objets irréductibles sont les $A_x$ $(x\in M)$, l'application $x\unsure{\mathrel{\vdash\joinrel\dashv}|\mapsto} A_x$ étant une bijection de $M$ sur l'ens des \deleted{\unsure{ses|son} \unsure{par|p}} objets irréductibles de $\Phi$.}

\deleted{\unsure{Dém|Dire|Bien} $A_x$ est irréductible, car si $A_x=\bigcup_{i\in I} A_i$}
\scanpage{72}\mspage{71}

Pour la suite, il faut faire attention si les axiomes vont passer aux sous-contrées. Introduisons quand $=\unsure{=}$ la définition \unsure{en forme|en jeu}

\uline{Définition 1} On appelle \uline{sous-contrée} \inserted{[\uline{pleine}]} de la contrée $\unsure{\deleted{\mathcal F}\mathcal C|\mathscr C|\mathbb C}=(\mathcal F,\leq,\unsure{\preceq|\leqslant})$, tout sous-ens. qui \unsure{satisfait|satisfont} les conditions

a) fermé pour $\leq$, et stable par sup \unsure{quelc.} pour $\leq$ [et à ce titre, s'en \unsure{correspond|correspondant}, \inserted{par $\mathcal F'\unsure{\mathrel{\vdash\joinrel\dashv}|\mapsto}\mathcal F'\cap M$,} aux parties fermées \inserted{$M'$} de $M$, par $M'\unsure{\mathrel{\mapstochar\relbar}|\mapsto}\mathcal F'=\{F\in\mathcal F\mid\overline F\subset M'\}$].

% slanted
\margin{Je ne sais s'il faut prendre des sup quelconques \unsure{(condition stricte)|(condition stricts)} \deleted{\unsure{et}} ou bien \unsure{de|des} sup finis seulement \unsure{(condition large)|(condition lge)}}

b) fermé pour $\preceq$ (ce qui impose une restriction aux parties envisagées $M'$ de $M$\unsure{\}|?}, \deleted{\unsure{à savoir|savoir} : si $X\in$} \inserted{qu'on va formuler} plus bas)

% Marginal continuation, written above C6, with a downward pointer.
\margin{\unsure{reporter|rajouter} \textcircled{c} Pour $F'\preceq F$ subd. \unsure{fixé|fixée}, l'application $G\mathrel{\vdash\joinrel\dashv}G'$, \deleted{\unsure{F}}sous-fig. de $F$ $\unsure{-|\longrightarrow}$ sous-fig. de $F'$ \deleted{\unsure{s}} commute aux Sup (pour $\leq$) \unsure{i.e.}}

\uline{C6} Soit

\diagram{\begin{tikzcd}[ampersand replacement=\&,column sep=1.4em,row sep=1em]
F' \arrow[r,phantom,"{\preceq}" description] \& F \arrow[d,phantom,"{\rotatebox{90}{$\leq$}}" description] \& \text{\unsure{de|dans|ds}}~\mathcal F \\
\& G \&
\end{tikzcd}}
% slanted
\margin{symétrique en $\leq,\preceq$}

\textcircled{a} Alors il existe un unique $H\in\mathcal F$ \unsure{avec|ou}

\diagram{\begin{tikzcd}[ampersand replacement=\&,column sep=1.4em,row sep=1em]
F' \arrow[r,phantom,"{\preceq}" description] \arrow[d,phantom,"{\rotatebox{90}{$\leq$}}" description] \& F \arrow[d,phantom,"{\rotatebox{90}{$\leq$}}" description] \\
G' \arrow[r,phantom,"{\preceq}" description] \& G
\end{tikzcd}}
\margin{et \textcircled{b} les \unsure{strates|str.tes|arêtes} de $G'$ sont les \deleted{\unsure{S}} \unsure{strates|str.tes|arêtes} $X'$ de $F'$ \unsure{tel|\deleted{tel}}les que $X'\ll\unsure{G|G_m|G_n}$. Ce complément \unsure{servira} plus bas}

On dit que $G'$ est la \uline{subd} \uline{de} $G$ induite par la subd $F'$ de $F$, ou que $G'$ est la \deleted{\unsure{sous-fig}}
% slanted
\margin{sur la sous-figure $G$ de $F$}
\scanpage{73}\mspage{72}

sous-figure de \deleted{$\unsure{G'|F'}$} la subdivision $F'$ de $F$, induite par la sous-figure $\deleted{\unsure{F|E}}G$ de $F$.
% slanted
\margin{\unsure{double} \unsure{Transitivité} de l'induction}

$\uline{C}7_0$ Pour toute figure $F$, sous-figure $G$, et subdivision $G'$ de $G$, il existe une subdivision $F'$ de $F$ qui \unsure{induise|induit} $\unsure{C'|G'}$.

(mais l'énoncé symétrique, où on échangerait les \unsure{rôle|rôles} de \deleted{\unsure{l|F}} $\leq$ et $\unsure{\preceq|\leqslant}$, est faux \unsure{\deleted{)}|?|\}|!} : il n'est pas vrai, pour $F'\preceq F$ \unsure{donnés|données}, que tte \unsure{sous figure|sous-figure} $G'$ de $F'$ \unsure{soit|sont|est} induite par une ss-figure $G$ de $\deleted{\unsure{G}}F$).

On \deleted{peut} \unsure{renforcera|renforce} par la suite $C7_0$ en \deleted{\unsure{s}}$C7$, en prenant un choix canonique de $F'$.

% Three diagonal cancellation strokes cover the draft below.
\deleted{\deleted{\uline{\unsure{Corollaire|Comparabilité} \unsure{1}}} Soient $F,\deleted{\unsure{A|F}}G$ deux figures.}

\deleted{\inserted{Considérons les} Conditions \deleted{\unsure{équivalentes}} :}

\deleted{a) $G$ est une sous-figure d'une subd. de $F$ :}
\[
\deleted{G\leq F'\preceq F}
\]
\deleted{b) $G$ est une subdivision d'une sous-figure de $F$ :}
\[
\deleted{G\preceq G'\leq F.}
\]
\deleted{On a alors b) $\Rightarrow$ a) (Mais pas \unsure{vrai|vice}\unsure{,} a $\Rightarrow$ b !!!)\unsure{.}}

\diagram{\begin{tikzcd}[ampersand replacement=\&,column sep=1.2em,row sep=1em]
\deleted{F} \arrow[r,phantom,"{\deleted{\unsure{\preceq\mkern-5mu{}^{\prime}|\preceq|\leqslant}}}" description] \arrow[d,phantom,"{\deleted{\unsure{\cup|U|\vee}}}" description] \& \deleted{F'} \arrow[d,dotted,->] \\
\deleted{G} \arrow[r,phantom,"{\deleted{\cdots}}" description] \& \deleted{G'}
\end{tikzcd}}

\deleted{\deleted{a) $\Rightarrow$ b) résulte de C6}}

\deleted{b) $\unsure{\Longrightarrow|\longrightarrow}$ a) résulte de $C7_0$}
